\lecture{17}{}{Fun with Squares}
\section{Squares}
\begin{note}
    The set $M_n(F)$ has special properties with respect to matrix multiplication:
    \begin{itemize}
        \item $M_n(F)$ is closed with respect to matrix multiplication (the product of two $n \times n$ matrices is always defined and is also $n \times n$)
        \item The matrix $I_n$ is the neutral element with respect to multiplication in $M_n(F)$
        \item The neutral element with respect to a binary operation is always unique; $I_n$ is the only matrix satisfying $I_n A = AI_n = A$ for all $A \in M_n(F)$
    \end{itemize}
\end{note}

\begin{note}
    However, this does not mean we cannot find $A, B \in M_n(F)$ with $I_n \neq A, B$ such that $AB = BA = A$.
\end{note}

\subsection{Commutativity}

\begin{note}
    In general, for matrices $A \in M_{m \times n}(F)$ and $B \in M_{p \times q}(F)$:
    \begin{itemize}
        \item $AB$ is not always defined
        \item Even when $AB$ is defined, $BA$ is not always defined
        \item Even when $AB$ and $BA$ are both defined, they might not have the same shape
        \item Even when $AB$ and $BA$ are both defined and have the same shape, they might not be equal
    \end{itemize}
\end{note}

\begin{definition}[Commuting Matrices]\label{def:commuting-matrices}
    Two square matrices $A$ and $B$ of the same order are said to \textbf{commute} if
    \[
    AB = BA
    \]
\end{definition}

\begin{note}
    A necessary condition for $AB = BA$ is that $A$ and $B$ must both be square and of the same order.
\end{note}

\begin{definition}[Scalar Matrix]\label{def:scalar-matrix}
    A matrix $A \in M_n(F)$ is called a \textbf{scalar matrix} if there exists $t \in F$ such that $A = tI_n$. That is, a scalar matrix is a diagonal matrix where all diagonal entries are equal.
\end{definition}

\begin{result}[Scalar Matrices Commute]\label{rslt:scalar-commute}
    Any scalar matrix commutes with any square matrix of the same order.
\end{result}

\begin{proof}
    Let $A, B \in M_n(F)$ and suppose $A$ is scalar. Then there exists $t \in F$ such that $A = tI_n$. Therefore:
    \begin{align*}
    AB &= (tI_n)B \\
    &= t(I_nB) && \text{(Proposition~\ref{prop:scalar-factoring}(i))} \\
    &= tB && \text{(Theorem~\ref{thm:identity-property}(i))} \\
    &= t(BI_n) && \text{(Theorem~\ref{thm:identity-property}(ii))} \\
    &= (tI_n)B && \text{(Proposition~\ref{prop:scalar-factoring}(ii))} \\
    &= B(tI_n) && \text{(Proposition~\ref{prop:scalar-factoring}(i))} \\
    &= BA
    \end{align*}
\end{proof}

\begin{theorem}[Diagonal Matrices Commute]\label{thm:diagonal-commute}
    Diagonal matrices of the same order commute.
\end{theorem}

\begin{proof}
    Let $A, B \in M_n(F)$ be diagonal. Then $AB$ and $BA$ are both defined and are square of order $n$. For all $i, j \in \{1, \ldots, n\}$:
    
    For $AB$:
    \begin{align*}
    [AB]_{ij} &= \sum_{k=1}^{n} [A]_{ik}[B]_{kj}
    \end{align*}
    
    Since $A$ is diagonal, $[A]_{ik} = 0$ whenever $i \neq k$. The only nonzero term in the sum occurs when $k = i$:
    \[
    [AB]_{ij} = [A]_{ii}[B]_{ij}
    \]
    
    Since $B$ is diagonal, $[B]_{ij} = 0$ when $i \neq j$. Therefore:
    \[
    [AB]_{ij} = [A]_{ii}[B]_{ii}\delta_{ij}
    \]
    
    Similarly, for $BA$:
    \begin{align*}
    [BA]_{ij} &= \sum_{k=1}^{n} [B]_{ik}[A]_{kj}
    \end{align*}
    
    Since $B$ is diagonal, the only nonzero term occurs when $k = i$:
    \[
    [BA]_{ij} = [B]_{ii}[A]_{ij}
    \]
    
    Since $A$ is diagonal, $[A]_{ij} = 0$ when $i \neq j$. Therefore:
    \[
    [BA]_{ij} = [B]_{ii}[A]_{ii}\delta_{ij} = [A]_{ii}[B]_{ii}\delta_{ij}
    \]
    
    Thus $[AB]_{ij} = [BA]_{ij}$ for all $i,j$, so $AB = BA$.
\end{proof}

\begin{note}
    We have found sufficient conditions for $AB = BA$:
    \begin{itemize}
        \item At least one of them is scalar
        \item Both of them are diagonal
    \end{itemize}
    However, there is no simple necessary and sufficient condition for commutativity.
\end{note}

\subsection{Matrix Powers}

\begin{definition}[Matrix Power]\label{def:matrix-power}
    Let $A \in M_n(F)$ be a square matrix and let $p \in \mathbb{N} \cup \{0\}$. The \textbf{matrix power} $A^p$ is defined recursively:
    \begin{itemize}
        \item If $p = 0$: $A^0 \coloneqq I_n$
        \item If $p > 0$: $A^p \coloneqq A^{p-1}A$
    \end{itemize}
\end{definition}

\begin{note}
    This definition gives us:
    \begin{itemize}
        \item $A^1 = A$
        \item $A^2 = AA$
        \item $A^3 = AAA$
        \item And so on
    \end{itemize}
\end{note}

\begin{proposition}[Powers Commute]\label{prop:powers-commute}
    If $B$ and $C$ are powers of the same matrix $A$, then they commute.
\end{proposition}

\begin{proof}
    Let $B = A^r$ and $C = A^s$ for some $r, s \in \mathbb{N} \cup \{0\}$. By associativity of matrix multiplication:
    \[
    A^r A^s = \underbrace{A \cdot A \cdots A}_{r+s \text{ times}} = A^s A^r
    \]
    Therefore $BC = CB$.
\end{proof}

\begin{eg}
    Consider the Fibonacci sequence defined by:
    \begin{itemize}
        \item $a_1 = 1$
        \item $a_2 = 1$
        \item $a_n = a_{n-1} + a_{n-2}$ for all $n > 2$
    \end{itemize}
    Matrix powers can be used to compute $a_{32}$ without computing all of $a_n$ for $n \in \{1, \ldots, 31\}$. (This technique uses the matrix $\begin{bmatrix} 1 & 1 \\ 1 & 0 \end{bmatrix}$ and its powers.)
\end{eg}